\begin{proposition}[Generic kernel is infinitesimal unipotent; Proved]
Let $X$ be a normal connected projective variety of dimension at least
two over an algebraically closed field $k$ of characteristic $p>0$.
For a specified very ample sheaf $L$, let $Y$ be the geometric generic
hyperplane section. Then the kernel
\[
 G=\ker\bigl(\mathbf{Pic}^{\tau}_{X/k}\otimes_k\overline{k(|L|)}
          \longrightarrow\mathbf{Pic}^{\tau}_{Y/\overline{k(|L|)}}\bigr)
\]
is finite, connected, and unipotent.
\end{proposition}
\begin{proof}
The finiteness proof in Proposition~\ref{a:generic} applies to normal
$X$: $\mathbf{Pic}^{\tau}$ is proper of finite type, and a
positive-dimensional reduced connected kernel would contain
prime-to-$p$ torsion, contrary to the fundamental-group surjection.

The argument excluding geometric points also works for normal $X$,
with one precision concerning differentials. Associate to a
$p$-torsion line bundle its rational logarithmic differential in
$\Omega^1_{k(X)/k}$, viewed as a regular form on $X_{\mathrm{sm}}$.
If it is zero, the local units describing a trivialization of the
$p$th tensor power have $p$th roots in $k(X)$. Normality implies that
each root and its inverse are regular on the relevant trivializing
open, so the roots trivialize the line bundle. Thus this rational
$d\log$ is injective on $\Pic(X)[p]$.

The incidence argument is carried out over $X_{\mathrm{sm}}$ and the
open locus of hyperplanes smooth at the specified point. Its relative
cotangent sheaf is locally free. A form defined over $k$ that vanishes
on the geometric generic hyperplane therefore vanishes throughout
this integral open incidence locus. At a smooth point where the form
is nonzero, very ampleness lets the tangent hyperplane vary over all
codimension-one subspaces. Some such subspace is not annihilated by
the form, a contradiction. Triviality of the restricted line bundle
implies vanishing of its relative logarithmic differential on the
hyperplane: its global units are constants. As in
Proposition~\ref{a:generic}, every finite-order geometric point of the
Picard scheme is defined over $k$. The argument excludes $p$-torsion,
and taking powers excludes all $p$-power torsion. Consequently $G$
has just its identity as a geometric point and is connected.

It remains to prove unipotence, which does not follow from
connectedness alone. Raynaud's Proposition~6.2.1, p.~50, and
Proposition~6.2.3(1), p.~51, give a natural isomorphism
\[
 \Hom(\mu_p,\mathbf{Pic}_{Z/\Omega})
       \simeq H^1_{\text{ét}}(Z,\ZZ/p)
\]
for a proper scheme $Z$ over an algebraically closed field $\Omega$
\cite{Raynaud1970PicardDuality}. Apply it to $X_\Omega$ and $Y$.
The surjection $\pi_1(Y)\to\pi_1(X_\Omega)$ makes restriction
injective on these cohomology groups. Hence $\Hom(\mu_p,G)=0$.

Here is an elementary finite-group argument completing the proof.
Let $E$ be the maximal étale quotient of the Cartier dual $G^D$.
Its dual $E^D$ is a closed subgroup of the finite connected scheme
$G$, and hence is connected. Over $\Omega$, $E$ is a constant finite
abelian group; connectedness of its diagonalizable dual implies
that it is a $p$-group. If $E\ne0$, a quotient $E\to\ZZ/p$ dualizes
to a subgroup $\mu_p\subset E^D\subset G$, a contradiction.
Thus $G^D$ is connected, so its finite coordinate algebra is local
with residue field $\Omega$.

Finite-dimensional representations of $G$ are modules over
$\mathcal O(G)^*=\mathcal O(G^D)$. That local algebra has only the
trivial simple module. Every such representation therefore has a
finite filtration with trivial one-dimensional factors. Applying
this to the faithful regular representation embeds $G$ in a group
of upper unitriangular matrices. This is unipotence.
\end{proof}

For normal $X$, the tangent-kernel computation also remains valid.
The geometric generic hyperplane is geometrically integral, so
$H^0(X_\Omega,\mathcal O)=H^0(Y,\mathcal O)=\Omega$; the sequence
$0\to L^{-1}\to\mathcal O_X\to\mathcal O_Y\to0$ identifies the tangent
space of $G$ with $H^1(X,L^{-1})\otimes_k\Omega$. A finite connected
group scheme is nontrivial exactly when its tangent space is nonzero,
by Nakayama applied to its local coordinate algebra.
